% ch12 — BSSN and CCZ4 Formulations
\chapter{BSSN and CCZ4 Formulations}
\label{ch:bssn-ccz4}

The raw ADM equations are only weakly hyperbolic and therefore unsuitable for
stable long-term numerical evolution. This chapter develops the two strongly
hyperbolic reformulations implemented in the codebase: the BSSN conformal
decomposition and the CCZ4 constraint-damping system.

\section{Why ADM Is Ill-Posed for Numerics}
\label{sec:bssn-adm-illposed}

We demonstrate the weak hyperbolicity of the ADM system by examining its
principal symbol and show how constraint-violating modes can grow without bound.

\section{BSSN Conformal Decomposition}
\label{sec:bssn-conformal}
\coderef{bssn}{rs}

We introduce the 24-field BSSN system: conformal factor $\phi$, conformal metric
$\tilde\gamma_{ij}$, trace-free extrinsic curvature $\tilde A_{ij}$, trace $K$,
and conformal connection functions $\tilde\Gamma^i$, and derive the evolution
equations.

\section{CCZ4 Constraint Damping}
\label{sec:ccz4-damping}
\coderef{ccz4}{rs}

We extend the Z4 system to its conformal-covariant form CCZ4, introducing the
28-field system with constraint-damping parameters $\kappa_1$ and $\kappa_2$
that exponentially suppress constraint violations.

\section{Gauge Conditions}
\label{sec:bssn-gauge}
\coderef{gauge}{rs}

We present the 1+log slicing condition for the lapse and the Gamma-driver
condition for the shift, explaining how these choices avoid coordinate
singularities and slice-stretching instabilities.

\section{Boundary Conditions}
\label{sec:bssn-boundary}

We discuss Sommerfeld (radiative) outgoing boundary conditions and the practical
issues of applying them on a finite computational domain.

\paragraph{Key References.}
The BSSN formulation follows \cite{Shibata1995,Baumgarte1999}; CCZ4 follows
\cite{Alic2012}; gauge conditions draw on \cite{Bona1995,Alcubierre2003}.
